\section{主要研究结论}
\label{sec:ch07_main_research_conclusions}

本文的仿真实现与指定案例表明，统一机器人模型、控制接口、任务接口和触觉观测合同，能够为多任务操作提供可复用的实验基础。候选转化为轨迹时，可达性筛选、桌面几何与局部接近约束、连续执行监测及尾段Strict检查承担不同层次的检验：几何质量不能替代IK可达性，可达性不能替代任务成功，任务成功也不能替代Strict示教质量。各层判断针对不同条件，当前证据仅支持已有模型、任务及案例范围，不能外推为通用抓取能力或实机有效性。

示教数据的可用性取决于语义和记录关系，而不只取决于文件能否读取或转换。策略观测与完整仿真状态、位置命令与IK目标需要区分，动作与观测的时序配对、数值有效性、版本来源和重复关系也应分别审查。指定修复与检查支持相应样本满足对应条件，历史命令不连续和状态字段缺失则说明格式合法不足以认定训练资格。文件Hash、Manifest和规模只能提供部分追踪依据，不能单独证明来源独立、全局无泄漏或物理真实性，也不能将局部修复扩大为全量生产认证。

跨域实验显示，局部指标改善尚不足以证明完整迁移。C27在原冻结Validation的11个实例中仅2个联合通过；$F_0$相对Original的SWD下降49.77\%，仅为指定合同下的边际改善。在第\ref{sec:ch06_multi_round_adversarial}节披露的修订后开发协议下，多轮对抗终点在$D_3$下FR达100\%，不能据此认定域不可辨识；多视角残差终点$F_{50}$虽Joint FR为100\%，分布距离却明显退化。P51因未满足既定Real接受要求对$D_{50}$判定NO-GO，该判别器不能作为当前实验的最终可靠评价器，该结论不外推$D_3$或其他方法。由此可见，接受率不等于分布相似，边际相似不保证输入输出对应，统计依赖不等于物理语义，梯度可达也不保证优化有效。评价器可靠性同样是迁移结论的前提。本文完成了统计伪监督、潜变量映射、共享参数向量场及多视角评价的系列探索，但现有证据尚不能证明完整触觉Sim-to-Real迁移成功。
